\documentclass[11pt]{article}
\usepackage[utf8]{inputenc}
\usepackage[T1]{fontenc}
\usepackage{lmodern}
\usepackage[margin=1.05in]{geometry}
\usepackage{amsmath,amssymb,amsthm,mathtools}
\usepackage{microtype,booktabs,array}
\usepackage{xcolor}
\usepackage{listings}
\usepackage[colorlinks=true,linkcolor=blue!45!black,citecolor=blue!45!black,urlcolor=blue!45!black]{hyperref}
\usepackage[nameinlink,capitalise,noabbrev]{cleveref}
\numberwithin{equation}{section}
\newtheorem{theorem}{Theorem}[section]
\newtheorem{proposition}[theorem]{Proposition}
\newtheorem{lemma}[theorem]{Lemma}
\newtheorem{corollary}[theorem]{Corollary}
\theoremstyle{definition}
\newtheorem{definition}[theorem]{Definition}
\theoremstyle{remark}
\newtheorem{remark}[theorem]{Remark}
\newcommand{\Q}{\mathbb Q}
\newcommand{\Z}{\mathbb Z}
\newcommand{\N}{\mathbb N}
\newcommand{\F}{\mathbb F}
\newcommand{\p}{\mathfrak p}
\newcommand{\PP}{\mathfrak P}
\DeclareMathOperator{\Gal}{Gal}
\DeclareMathOperator{\Frob}{Frob}

\lstdefinestyle{lean}{
  basicstyle=\ttfamily\footnotesize,
  columns=fullflexible,
  keepspaces=true,
  breaklines=true,
  frame=single,
  framerule=0.3pt,
  rulecolor=\color{black!40},
  xleftmargin=0.5em,
  morekeywords={def,theorem,fun,if,then,else,let,mut,for,in,do,return,by,Nat,Prop,Bool},
  keywordstyle=\bfseries,
  extendedchars=true,
  inputencoding=utf8,
  literate=
    {∀}{{$\forall$}}1 {∃}{{$\exists$}}1 {≤}{{$\le$}}1 {≥}{{$\ge$}}1
    {→}{{$\to$}}1 {∣}{{$\mid$}}1 {¬}{{$\neg$}}1 {∧}{{$\wedge$}}1
    {∨}{{$\vee$}}1 {≠}{{$\ne$}}1 {Λ}{{$\Lambda$}}1 {≡}{{$\equiv$}}1
    {…}{{$\ldots$}}1 {·}{{$\cdot$}}1 {ℕ}{{$\mathbb N$}}1
}

\title{Explicit lower bounds for runs of three consecutive\\ quintic and septic residues}
\author{Claude (AI model by Anthropic)\\ \small run and submitted by Aleksandr Bersenev}
\date{October 8, 2026\\[3pt]\footnotesize Produced autonomously by Claude (Opus 5.5) agents in Claude Code; run and submitted by Aleksandr Bersenev.\\ \footnotesize Not an official Anthropic publication. Not peer reviewed.}
\hypersetup{pdftitle={Explicit lower bounds for runs of three consecutive quintic and septic residues},pdfauthor={Claude (AI model by Anthropic); submitted by Aleksandr Bersenev}}

\begin{document}
\maketitle

\begin{abstract}
For a prime $p$ let $r(k,3,p)$ be the least $r\ge1$ such that $r,r+1,r+2$ are all $k$-th power
residues modulo $p$, and $\Lambda(k,3)=\limsup_{p\to\infty}r(k,3,p)$.
Lehmer, Lehmer, Mills and Selfridge proved $\Lambda(3,3)=23532$; whether $\Lambda(k,3)$ is finite for
odd $k\ge5$ is open (Erd\H{o}s Problem \#436).
We give explicit assignments of quintic (resp.\ septic) characters to the primes and prove, with a
Lean~4 certificate (one \texttt{native\_decide} evaluation, everything else kernel-checked), that no
completely additive function realising them has three consecutive zeros
below $10^8$ (resp.\ $10^7$). Combined with the classical Kummer--Chebotarev realisability argument,
which we prove in the paper but did not formalise, this gives
$\Lambda(5,3)\ge 99\,999\,999$ and $\Lambda(7,3)\ge 9\,999\,999$.
Computations outside Lean show that the same assignments have their first triples of
consecutive zeros at $313\,714\,024$ and $180\,317\,500$, so
$\Lambda(5,3)\ge313\,714\,024$ and $\Lambda(7,3)\ge180\,317\,500$ (computational, not formally
certified). The bounds previously available were the trivial $\Lambda(5,3)\ge\Lambda(5,2)=7888$ and
$\Lambda(7,3)\ge\Lambda(7,2)=1\,649\,375$. Whether $\Lambda(5,3)$ and $\Lambda(7,3)$ are finite remains open.
\end{abstract}

\tableofcontents

%======================================================================
\section{Introduction}

Let $k,m\ge2$. By a theorem of Brauer \cite{Br28}, every sufficiently large prime $p$ has $m$ consecutive
positive integers $r,r+1,\dots,r+m-1$ that are all $k$-th power residues modulo $p$; let
$r(k,m,p)$ denote the least such $r$ and put
\[
  \Lambda(k,m)=\limsup_{p\to\infty} r(k,m,p)\in\N\cup\{\infty\}.
\]
D.~H.~Lehmer and E.~Lehmer \cite{LeLe62} started the study of $\Lambda$: they observed
$\Lambda(2,2)=9$ (if neither $2$ nor $5$ is a quadratic residue then $10$ is), and proved
$\Lambda(2,3)=\infty$ and $\Lambda(k,4)=\infty$ for $k\le 1048909$; more generally $\Lambda(k,m)=\infty$
whenever $k$ is even and $m\ge3$ \cite{LeLe62,LLM63}, and Graham \cite{Gr64} showed
$\Lambda(k,m)=\infty$ for all $k$ and all $m\ge4$.
The finite values known are
\begin{gather*}
\Lambda(3,2)=77\ \text{\cite{Du65}},\qquad \Lambda(4,2)=1224\ \text{\cite{BiMi63}},\qquad
\Lambda(5,2)=7888,\quad \Lambda(6,2)=202124\ \text{\cite{LLM63}},\\
\Lambda(7,2)=1649375\ \text{\cite{BLL64}},
\end{gather*}
and, for runs of length three, the ``machine proof''
\[
  \Lambda(3,3)=23532 \qquad\text{(Lehmer, Lehmer, Mills and Selfridge \cite{LLMS62}).}
\]
Hildebrand proved that $\Lambda(k,2)<\infty$ for every prime $k$ \cite{Hi86} and then for every $k$
\cite{Hi91} (a combinatorial proof via a generalised Schur theorem is given by Dietzel \cite{Di13}).
Brillhart, Lehmer and Lehmer \cite{BLL64} also proved the general lower bound $\Lambda(k,2)>5^k$ for
every odd $k\ge5$ (and for $k$ twice an odd prime $\equiv3\pmod4$).

Erd\H{o}s and Graham \cite[p.~87]{ErGr80} asked whether $\Lambda(k,3)$ is finite for every odd $k$ and how
large $\Lambda(k,2)$ and $\Lambda(k,3)$ are; this is Problem~\#436 on the Erd\H{o}s problems website
\cite{EP436}, which (page last edited 25 October 2025, consulted October 8, 2026) records as the
remaining questions ``whether $\Lambda(k,3)$ is finite for all odd $k\ge5$, and the growth rate of
$\Lambda(k,2)$ and $\Lambda(k,3)$ as functions of $k$''. The first open case is $k=5$.

\paragraph{What is new.}
Every run of three residues begins with a pair, so $r(k,3,p)\ge r(k,2,p)$ and trivially
$\Lambda(5,3)\ge\Lambda(5,2)=7888$, $\Lambda(7,3)\ge\Lambda(7,2)=1\,649\,375$ (and $\Lambda(k,3)\ge\Lambda(k,2)>5^k$
for every odd $k\ge5$ by \cite{BLL64}; the $\Omega$-assignment of \cref{sec:comp} gives
$\Lambda(k,3)\ge3^k-2$). To our knowledge, as of October 8, 2026, no bound specific to runs of
\emph{three} quintic or septic residues, i.e.\ going beyond these trivial consequences of the
pair results, has been published. We give the first such bounds:

\begin{theorem}[main result]\label{thm:main}
\[
  \Lambda(5,3)\ \ge\ 99\,999\,999 \qquad\text{and}\qquad \Lambda(7,3)\ \ge\ 9\,999\,999 .
\]
More precisely, for the explicit prime assignments $g_5,g_7$ of \cref{sec:assign}, there are infinitely
many primes $P\equiv1\pmod5$ (resp.\ $P\equiv1\pmod7$) with $r(5,3,P)\ge 10^8-1$ (resp.\
$r(7,3,P)\ge10^7-1$).
\end{theorem}

\paragraph{Logical structure.} The proof has two parts of different status.
\begin{enumerate}
\item A purely combinatorial statement about completely additive functions with prescribed prime
values (Theorem~\ref{thm:core}). It is \emph{certified in Lean~4} (core library only, no Mathlib). The
certificate uses \texttt{native\_decide} exactly once, to evaluate one large finite check on a table
of $10^8+1$ bytes; everything else is checked by the Lean kernel (\cref{sec:lean}).
\item The passage from completely additive functions to $k$-th power residues modulo infinitely many
primes $P$: a classical consequence of Kummer theory and the Chebotarev density theorem
(Lemma~\ref{lem:real}). We give a complete proof but have \emph{not} formalised it.
\end{enumerate}
In addition we report stronger bounds that are \emph{computational only} (\cref{sec:comp}):
\[
  \Lambda(5,3)\ \ge\ 313\,714\,024,\qquad \Lambda(7,3)\ \ge\ 180\,317\,500 .
\]

\paragraph{What is not shown.} We do not know whether $\Lambda(5,3)$ or $\Lambda(7,3)$ is finite.
Our searches found no indication of a threshold beyond which every assignment of characters
produces a zero triple (\cref{sec:open}); if $\Lambda(5,3)$ is finite, it exceeds $3\cdot10^8$, about
$13\,000$ times $\Lambda(3,3)$.

%======================================================================
\section{Reduction to completely additive functions}\label{sec:reduction}

Throughout, $k$ is an odd prime. For a prime $P\equiv1\pmod k$ the $k$-th power residues modulo $P$ form
the subgroup of index $k$ of the cyclic group $\F_P^\times$. Fix a primitive $k$-th root of unity
$\zeta=\zeta_k$ in $K=\Q(\zeta)$.

\begin{definition}
An \emph{assignment} is a map $g$ from the primes to $\Z/k\Z$. A function $f:\N\to\N$ is
\emph{completely additive mod $k$ on $[1,L]$} if $f(ab)\equiv f(a)+f(b)\pmod k$ whenever $a,b\ge1$ and
$ab\le L$; it \emph{realises $g$ on $[1,L]$} if moreover $f(q)\equiv g(q)\pmod k$ for every prime
$q\le L$. (Such $f$ are unique mod $k$ on $[1,L]$; see the proof of Theorem~\ref{thm:core}.)
\end{definition}

\begin{lemma}[realisability]\label{lem:real}
Let $k$ be an odd prime, $q_1<\dots<q_t$ distinct primes ($k$ itself may be among them) and
$c\in(\Z/k\Z)^t$. There are infinitely many primes $P\equiv1\pmod k$ for which there is a map
$f_P:\{n\in\Z: P\nmid n\}\to\Z/k\Z$ such that
\begin{enumerate}
\item $f_P(mn)=f_P(m)+f_P(n)$;
\item $f_P(q_i)=c_i$ for $i=1,\dots,t$;
\item $n$ is a $k$-th power residue modulo $P$ if and only if $f_P(n)=0$.
\end{enumerate}
\end{lemma}

\begin{proof}
\emph{Independence.} Let $K=\Q(\zeta)$, so $[K:\Q]=k-1$ and $\mu_k\subset K$. Suppose
$x=\prod_i q_i^{a_i}=\alpha^k$ with $\alpha\in K^\times$ and $0\le a_i<k$. Taking the norm
$N_{K/\Q}$, and using that $x\in\Q$,
\[
  \prod_i q_i^{(k-1)a_i}=N_{K/\Q}(x)=N_{K/\Q}(\alpha)^k\in(\Q^\times)^k .
\]
Comparing $q_i$-adic valuations gives $k\mid(k-1)a_i$, hence $k\mid a_i$ and $a_i=0$. So the classes of
$q_1,\dots,q_t$ are $\F_k$-linearly independent in $K^\times/(K^\times)^k$.

\emph{Kummer theory.} Let $\Delta\subseteq K^\times$ be generated by $(K^\times)^k$ and $q_1,\dots,q_t$,
and $M=K(\Delta^{1/k})=K(q_1^{1/k},\dots,q_t^{1/k})$ for fixed choices of the roots. Since
$\mu_k\subset K$, Kummer theory gives $[M:K]=|\Delta/(K^\times)^k|=k^t$, and the map
\[
  \Gal(M/K)\to\mu_k^t,\qquad \sigma\mapsto\bigl(\sigma(q_i^{1/k})/q_i^{1/k}\bigr)_{i}
\]
is an injective homomorphism between groups of order $k^t$, hence an isomorphism. Let $\sigma_c$ be
the element with $\sigma_c(q_i^{1/k})=\zeta^{c_i}q_i^{1/k}$ for all $i$.

\emph{Chebotarev.} $M/K$ is abelian and unramified outside the primes dividing $kq_1\cdots q_t$. By the
Chebotarev density theorem the primes $\p$ of $K$, unramified in $M$, with Frobenius
$\Frob_\p=\sigma_c$ have Dirichlet density $k^{-t}>0$. Primes $\p$ of residue degree $\ge2$ over $\Q$
have density $0$ (at most $(k-1)\pi(\sqrt x)$ of them have norm $\le x$), so infinitely many $\p$ in
this class have degree one: $N\p=P$ is a rational prime, $P\nmid kq_1\cdots q_t$, and $P$ splits in
$K$, i.e.\ $P\equiv1\pmod k$. Fix such a $\p$; then $O_K/\p\cong\F_P$.

\emph{Power residue symbols.} Let $\PP$ be a prime of $M$ above $\p$ and $\beta=q_i^{1/k}$. By definition
of the Frobenius, $\sigma_c(\beta)\equiv\beta^{P}\pmod\PP$, so
$\zeta^{c_i}\beta\equiv\beta^P$, i.e.\ $q_i^{(P-1)/k}=\beta^{P-1}\equiv\zeta^{c_i}\pmod\PP$, since $\beta$ is
a unit at $\PP$. Both sides lie in $O_K$, so the congruence holds modulo $\PP\cap O_K=\p$.

\emph{Definition of $f_P$.} Since $P\nmid k$, the polynomial $X^k-1$ is separable modulo $P$, so the $k$
elements of $\mu_k$ remain distinct modulo $\p$ and reduce bijectively onto the $k$-th roots of unity
of $\F_P$ (there are exactly $k$ of them, as $k\mid P-1$). For $P\nmid n$ the element
$n^{(P-1)/k}\bmod\p$ satisfies $(n^{(P-1)/k})^k=n^{P-1}\equiv1$, so there is a unique
$f_P(n)\in\Z/k\Z$ with $n^{(P-1)/k}\equiv\zeta^{f_P(n)}\pmod\p$. Property~(1) follows from
multiplicativity of $n\mapsto n^{(P-1)/k}$ and injectivity of reduction on $\mu_k$; property~(2) is the
previous paragraph. For~(3): $\F_P^\times$ is cyclic of order $P-1$ and $k\mid P-1$, so $n$ is a $k$-th
power modulo $P$ iff $n^{(P-1)/k}\equiv1\pmod P$ (Euler's criterion), iff $f_P(n)=0$.
\end{proof}

\begin{remark}
Lemma~\ref{lem:real} is classical and not new. For odd prime $k$ and primes $q_i\ne k$ it is the
Lemma on p.~104 of Lehmer and Lehmer \cite{LeLe62}, stated there as a special case of a theorem of
Kummer (Hilbert's \emph{Zahlbericht} \cite{Hi97}); Mills \cite{Mi63} proved the version for arbitrary $k$,
which is the statement used by Lehmer, Lehmer and Mills: ``every vector $[r_1,\dots,r_t]$ \dots{} that
satisfies certain parity conditions is the $S$-vector of infinitely many primes'' \cite[p.~174]{LLM63};
for odd prime $k$ there are no parity conditions. Graham \cite{Gr64} uses it in exactly the form above
(``a special case of a result of Kummer''). We include a proof only because the Lean statement
requires the precise hypotheses, and because we need $k$ itself among the $q_i$ ($g_5(5)=0$,
$g_7(7)=1$), a case excluded in \cite{LeLe62}. All known values of $\Lambda$ were obtained through this reduction.
\end{remark}

\begin{corollary}\label{cor:transfer}
Let $k$ be an odd prime, $g$ an assignment and $L\ge3$. Suppose that every $f:\N\to\N$ that is completely
additive mod $k$ on $[1,L]$ and realises $g$ on $[1,L]$ has no $r$ with $1\le r$, $r+2\le L$ and
$f(r)\equiv f(r+1)\equiv f(r+2)\equiv0\pmod k$. Then $r(k,3,P)\ge L-1$ for infinitely many primes
$P$; in particular $\Lambda(k,3)\ge L-1$.
\end{corollary}

\begin{proof}
Apply Lemma~\ref{lem:real} to the primes $q\le L$ and $c_q=g(q)$, and take any of the infinitely many
primes $P>L$ it provides. Define $f(n)\in\{0,\dots,k-1\}$ as the representative of $f_P(n)$ for
$P\nmid n$ and $f(n)=0$ otherwise. Every $n\le L$ is prime to $P$, so $f$ is completely additive mod $k$
on $[1,L]$ and realises $g$ there. By hypothesis no $r\le L-2$ has $r,r+1,r+2$ all $k$-th power
residues modulo $P$, so $r(k,3,P)\ge L-1$.
\end{proof}

%======================================================================
\section{The assignments}\label{sec:assign}

\begin{definition}\label{def:g}
The assignment $g_5$ (values in $\Z/5\Z$) is $g_5(q)=1$ for every prime $q$ except the $25$ primes
in \cref{tab:g5}. The assignment $g_7$ (values in $\Z/7\Z$) is $g_7(q)=1$ for every prime $q$ except
$g_7(2)=g_7(47)=g_7(163)=0$.
\end{definition}

\begin{table}[h]
\centering\small
\setlength{\tabcolsep}{4pt}
\begin{tabular}{@{}l*{13}{r}@{}}
\toprule
$q$      & 2 & 5 & 7 & 11 & 13 & 17 & 23 & 31 & 73 & 347 & 409 & 617 & 1033 \\
$g_5(q)$ & 0 & 0 & 0 & 0  & 0  & 0  & 0  & 0  & 0  & 2   & 2   & 2   & 2 \\
\midrule
$q$      & 1373 & 2053 & 2089 & 2099 & 2437 & 2621 & 3907 & 4523 & 7057 & 12917 & 21407 & 49613 & \\
$g_5(q)$ & 0    & 0    & 2    & 2    & 2    & 0    & 2    & 0    & 2    & 2     & 3     & 3 & \\
\bottomrule
\end{tabular}
\caption{The exceptional primes of $g_5$; all other primes have $g_5(q)=1$.}\label{tab:g5}
\end{table}

For an assignment $g$ write $\chi_g(n)=\sum_{q}v_q(n)\,g(q)\in\Z/k\Z$ for its completely additive
extension. With $g\equiv1$ (the ``$\Omega$-assignment'') $\chi_g(n)=\Omega(n)\bmod k$; its first triple
of consecutive zeros is at $15470$ for $k=5$ and at $3145310$ for $k=7$. The assignments $g_5,g_7$ were
found by a local search (\cref{sec:comp}) that starts from the $\Omega$-assignment and changes the values
of a few primes until no zero triple remains below a target bound.

%======================================================================
\section{The Lean-certified core}\label{sec:lean}

\begin{theorem}[certified in Lean]\label{thm:core}
\leavevmode
\begin{enumerate}
\item Let $L=10^8$. Every $f:\N\to\N$ that is completely additive mod $5$ on $[1,L]$ and realises
$g_5$ on $[1,L]$ satisfies: there is no $r$ with $1\le r$, $r+2\le L$ and
$f(r)\equiv f(r+1)\equiv f(r+2)\equiv0\pmod5$.
\item The same holds for $L=10^7$, modulus $7$ and $g_7$.
\end{enumerate}
\end{theorem}

\Cref{thm:main} follows from Theorem~\ref{thm:core} and Corollary~\ref{cor:transfer}.

\subsection{The Lean statement}
The file \texttt{lean/Main.lean} (Lean 4.34.1, core library only, no imports) contains, verbatim
(docstrings omitted):

\begin{lstlisting}[style=lean]
def L : Nat := 100000000
def IsPrime (p : Nat) : Prop := 2 ≤ p ∧ ∀ d, d ∣ p → d = 1 ∨ d = p
def g (p : Nat) : Nat :=
  if 49613 < p then 1 else
  if p = 2 then 0 else if p = 5 then 0 else if p = 7 then 0 else if p = 11 then 0 else
  if p = 13 then 0 else if p = 17 then 0 else if p = 23 then 0 else if p = 31 then 0 else
  if p = 73 then 0 else if p = 347 then 2 else if p = 409 then 2 else if p = 617 then 2 else
  if p = 1033 then 2 else if p = 1373 then 0 else if p = 2053 then 0 else if p = 2089 then 2 else
  if p = 2099 then 2 else if p = 2437 then 2 else if p = 2621 then 0 else if p = 3907 then 2 else
  if p = 4523 then 0 else if p = 7057 then 2 else if p = 12917 then 2 else if p = 21407 then 3 else
  if p = 49613 then 3 else 1

theorem no_three_consecutive_quintic_residues (f : Nat → Nat)
    (hmul : ∀ a b, 1 ≤ a → 1 ≤ b → a * b ≤ L → f (a * b) % 5 = (f a + f b) % 5)
    (hprime : ∀ p, p ≤ L → IsPrime p → f p % 5 = g p % 5) :
    ∀ r, 1 ≤ r → r + 2 ≤ L → ¬ (f r % 5 = 0 ∧ f (r + 1) % 5 = 0 ∧ f (r + 2) % 5 = 0)
\end{lstlisting}

\texttt{lean/Main7.lean} is the same file with modulus $7$, \texttt{L := 10000000}, \texttt{B := 3163} and
\begin{lstlisting}[style=lean]
def g (p : Nat) : Nat :=
  if 163 < p then 1 else
  if p = 2 then 0 else if p = 47 then 0 else if p = 163 then 0 else 1
\end{lstlisting}
and main theorem \texttt{no\_three\_consecutive\_septic\_residues}.

\subsection{Proof architecture}
The proof introduces a table and a bound that do \emph{not} occur in the statement:
\begin{itemize}
\item \texttt{T : ByteArray := buildT L}, an array of $L+1$ bytes in which a sieve stores
$T[n]=\sum_{p^j\mid n}g(p)\bmod 5$, i.e.\ $\chi_{g}(n)$. The sieve is \emph{not trusted}: only the
properties below, checked on the computed array, are used.
\item \texttt{B := 10000}, with $L<(B+1)^2$.
\end{itemize}
A single Boolean \texttt{bigCheck} verifies, for the computed table:
\begin{enumerate}
\item[(i)] $T[1]\equiv0$ and $L<(B+1)^2$;
\item[(ii)] for every $p\le B$ that has no divisor $e$ with $2\le e$, $e^2\le p$ (a trial-division
primality test) and every $2\le m\le L/p$: $T[pm]\equiv T[p]+T[m]$;
\item[(iii)] for every $n\in[2,L]$: $T[n]\equiv g(n)$, or $n$ has a divisor $e$ with $2\le e$, $e^2\le n$;
\item[(iv)] there is no $r\in[1,L-2]$ with $T[r]\equiv T[r+1]\equiv T[r+2]\equiv0$.
\end{enumerate}
The theorem \texttt{bigCheck\_true : bigCheck = true} is proved by \texttt{native\_decide}. All remaining
steps are kernel-checked:
\begin{itemize}
\item soundness lemmas turn \texttt{bigCheck = true} into the Propositions (i)--(iv)
(\texttt{loopAll\_sound}, \texttt{pairLoop\_sound}, \texttt{hsd\_sound}, \texttt{quickComp\_sound});
\item \texttt{agree}: by strong induction, $f(n)\equiv T[n]$ for all $n\in[1,L]$. For $n=1$,
$f(1)\equiv f(1)+f(1)$ gives $f(1)\equiv0\equiv T[1]$. If $n$ is prime, (iii) applies with its second
alternative impossible (a divisor $e$ with $2\le e$, $e^2\le n$ would be a proper divisor), so
$T[n]\equiv g(n)\equiv f(n)$. If $n\ge2$ is not prime, write $n=ab$ with $2\le a\le b$
(\texttt{split\_composite}), take a prime $p\mid a$ (\texttt{exists\_prime\_dvd}); then $p^2\le a^2\le n\le L$,
so $p\le B$ by (i), $p$ passes the primality test in (ii), and with $n=pm$ we have $m\ge p\ge2$, whence
$f(n)\equiv f(p)+f(m)\equiv T[p]+T[m]\equiv T[n]$ by \texttt{hmul}, the induction hypothesis and (ii);
\item the main theorem then follows from (iv).
\end{itemize}

\subsection{Axioms, the use of \texttt{native\_decide}, and runs}
\texttt{\#print axioms} reports, for both main theorems,
\begin{center}\small
\texttt{[propext, Classical.choice, Quot.sound, bigCheck\_true.\_native.native\_decide.ax\_1\_1]}
\end{center}
and for \texttt{bigCheck\_true} itself
\begin{center}\small
\texttt{[propext, Quot.sound, bigCheck\_true.\_native.native\_decide.ax\_1\_1]}.
\end{center}
There is no \texttt{sorry}, \texttt{admit} or user-declared axiom. The last axiom expresses trust in Lean's
evaluator for the single closed Boolean term \texttt{bigCheck}: the trusted base of Theorem~\ref{thm:core}
is the Lean kernel \emph{plus} the compiler/IR interpreter invoked by \texttt{native\_decide}
(\texttt{Lean.ofReduceBool}). A reader who does not extend trust to the latter should regard
Theorem~\ref{thm:core} as computer-checked rather than kernel-checked; the two independent C programs
of \cref{sec:comp} confirm the same table-level facts.

\paragraph{Why not the kernel.} \texttt{bigCheck} builds an array of $10^8+1$ bytes by a sieve and then
performs of the order of $10^9$ table look-ups (checks (ii)--(iv)). The Lean kernel evaluates terms by
reduction, without a constant-time array representation; even the interpreter needed about 82 minutes
for $L=10^8$. A kernel reduction of this size is far out of reach, so the check is delegated to
\texttt{native\_decide}, and only to this one closed term.

\paragraph{Runs.} Lean 4.34.1 on a shared 14-core machine:
\texttt{Main.lean} ($k=5$, $L=10^8$): wall time $4940$\,s, 725\,MB;
\texttt{Main7.lean} ($k=7$, $L=10^7$): $165$\,s, 540\,MB. Both were run with plain
\texttt{lean <file>} (no project file, no Mathlib) under a timing/memory-limit wrapper;
\texttt{Main7.lean} with the extra flag \texttt{-j 2}. An independent re-run by the reviewer with the
same toolchain reproduced both outputs (see \texttt{README.md}).
Two control files (in \texttt{lean/}) show that the check is not vacuous: with $g\equiv1$ the same
\texttt{bigCheck} evaluates to \texttt{false} for $L=15500$ (zero triple $15470$--$15472$) and to
\texttt{true} for $L=15469$. An earlier version of the checker, organised differently but with the same
proof plan, certified the $k=5$ statement for $L=10^7$.

\paragraph{Statement audit.} The hypotheses quantify over \emph{all} $f:\N\to\N$; \texttt{hmul} is
complete additivity mod $5$ on $[1,L]$ (all $a,b\ge1$ with $ab\le L$), and \texttt{hprime} prescribes
$f$ on every prime $\le L$ via \texttt{IsPrime}, the textbook definition. These are exactly the
properties of the index functions $f_P$ of Lemma~\ref{lem:real} for $P>L$. The function \texttt{g} is a
plain if-chain whose values agree with Definition~\ref{def:g}; the leading test \texttt{49613 < p} only
short-cuts the chain (no exceptional prime exceeds $49613$). The conclusion covers every window
$r,r+1,r+2$ contained in $[1,L]$, so a first zero triple must start at $r\ge L-1$, as used in
Corollary~\ref{cor:transfer}. The table \texttt{T}, the bound \texttt{B} and the sieve do not appear in the
statement.

%======================================================================
\section{Computational results (not formally certified)}\label{sec:comp}

\begin{proposition}[computation]\label{prop:comp}
The completely additive extension $\chi_{g_5}$ has no triple of consecutive zeros $r,r+1,r+2$ with
$r+2\le313\,714\,025$, and
\[
\begin{array}{rcl@{\qquad}l}
313714024 &=& 2^3\cdot13^4\cdot1373, & \chi_{g_5}=3\cdot0+4\cdot0+0=0,\\
313714025 &=& 5^2\cdot29^2\cdot43\cdot347, & \chi_{g_5}=0+2+1+2\equiv0,\\
313714026 &=& 2\cdot3^3\cdot997\cdot5827, & \chi_{g_5}=0+3+1+1\equiv0 \pmod5.
\end{array}
\]
Likewise $\chi_{g_7}$ has no zero triple with $r+2\le180\,317\,501$, and its first zero triple is
$180317500=2^2\cdot5^4\cdot11\cdot79\cdot83$, $180317501=7^4\cdot13\cdot53\cdot109$,
$180317502=2\cdot3^4\cdot37\cdot67\cdot449$.
\end{proposition}

By Corollary~\ref{cor:transfer} (with $L=313\,714\,025$, resp.\ $L=180\,317\,501$), Proposition~\ref{prop:comp} gives $\Lambda(5,3)\ge313\,714\,024$ and
$\Lambda(7,3)\ge180\,317\,500$; moreover $r(5,3,P)=313\,714\,024$ for infinitely many $P$, namely the primes
of Lemma~\ref{lem:real} for the primes $q\le313\,714\,026$.

\paragraph{How this was checked.} Two independent C programs compute $\chi_g$ and scan for zero triples:
one uses a smallest-prime-factor table and the recursion $\chi(n)=g(p)+\chi(n/p)$; the other factors each
segment of $2^{24}$ integers by trial division by all primes up to $\sqrt{L}$. Both find the same
first triples. The three factorisations above were also checked with SymPy and can be verified by hand.
These computations are not formally certified. \texttt{Main.lean} with \texttt{L := 313714025} and
\texttt{B := 17713} (so that $(B+1)^2>L$) would certify the bound $\Lambda(5,3)\ge313\,714\,024$, but
we did not run it: at the measured interpreter speed it would need several hours.

\paragraph{Search.} The assignments were found by a WalkSAT-type local search over the values of the
primes, with incremental bookkeeping of zero triples: a zero window is picked at random and the value of
one of the primes dividing its members is changed so as to minimise the number of zero windows. Starting
from the $\Omega$-assignment, 4 and 36 moves removed all zero triples up to $10^6$ and $10^7$ ($k=5$);
61 further moves extended a $10^7$ solution to $10^8$; for $k=7$, three moves sufficed up to $10^8$.
An exact depth-first search with conflict-directed backjumping over the same model reproduces the known
values $\Lambda(2,2)=9$, $\Lambda(3,2)=77$, $\Lambda(4,2)=1224$ and $\Lambda(5,2)=7888$ as the boundary
between solvable and unsolvable bounds, which checks the model (including Lemma~\ref{lem:real}) against the
literature for $k=2,3,5$. (For $k=4$ this is only a plausibility check: Lemma~\ref{lem:real} is stated
for odd prime $k$, and the symmetry reduction used by the search---scaling all values by a unit of
$\Z/k\Z$, so that the first non-zero prime value may be taken to be $1$---is transitive on the non-zero
values only for prime $k$; for $k=4$ the search therefore omits assignments whose first non-zero value
is $2$, and its ``unsolvable'' verdict at $1225$ is not a complete search.)

\paragraph{Side data.} The $\Omega$-assignment alone gives $\Lambda(11,2)\ge6\,082\,047$ (first zero
pair $6082047=3^{10}\cdot103$, $6082048=2^9\cdot7\cdot1697$) and
$\Lambda(13,2)\ge10^8$ (no two consecutive integers below $10^8$ have $\Omega\equiv0\pmod{13}$), again by
computation only. These two bounds are \emph{weaker} than what has been known since 1964: Brillhart,
Lehmer and Lehmer \cite[\S\S1, 6]{BLL64} proved $\Lambda(k,2)>5^k$ for every odd $k\ge5$ by an
elementary argument (the assignment $R(3)=0$, $R(q)=1$ for $q\ne3$, and factorisations of $3^n\pm1$),
so $\Lambda(11,2)>48\,828\,125$ and $\Lambda(13,2)>1.2\cdot10^9$; they remark that $7\cdot5^{k-1}$ can be
obtained similarly, and they give $\Lambda(8,2)\ge1\,200\,744$, $\Lambda(9,2)>10^7$,
$\Lambda(10,2)\ge22\,458\,303$. Elementarily, for every odd prime $k$ the
$\Omega$-assignment shows $\Lambda(k,2)\ge3^k-1$ and $\Lambda(k,3)\ge3^k-2$: every member $n\ge2$ of a zero
run has $\Omega(n)\ge k$, and one member is odd, hence at least $3^k$. We claim no novelty for these remarks.

%======================================================================
\section{Finiteness remains open}\label{sec:open}

We stress that Theorem~\ref{thm:main} and Proposition~\ref{prop:comp} are lower bounds; they do not decide whether
$\Lambda(5,3)$ or $\Lambda(7,3)$ is finite.

By Lemma~\ref{lem:real}, $\Lambda(k,3)$ is finite if and only if for some $L$ every assignment $g$ has a zero
triple of $\chi_g$ in $[1,L]$ (if not, Corollary~\ref{cor:transfer} gives $\Lambda(k,3)\ge L-1$ for every $L$;
conversely every prime $P\equiv1\pmod k$ induces an assignment). We found no sign of such an $L$: the exact search slows down near $1.4\cdot10^5$ without proving
unsolvability, while the local search finds solutions far beyond with very few changes.

The method behind $\Lambda(3,3)=23532$ \cite{LLMS62} and the values of $\Lambda(k,2)$ restricts to the
first $t$ primes and uses runs of consecutive $p_t$-smooth integers, whose zero conditions are linear
forms on $(\Z/k\Z)^t$. By the analogue for triples of the heuristic count in \cite[p.~174]{LLM63}, covering $(\Z/5\Z)^{25}$ by the
codimension-three subspaces coming from triples would need about
$-t\log5/\log(1-5^{-3})\approx5000$ triples of consecutive $97$-smooth integers, while there are only
$787$ such triples below $10^{12}$ (our count). So the approach of 1962 does not extend in an obvious way.
Nor can periodic constructions prove $\Lambda(k,3)=\infty$ for odd $k$: if $\chi$ agrees on integers
prime to $M$ with a character $\psi$ of $(\Z/M\Z)^\times$ of order dividing $k$ (written additively,
with values in $\Z/k\Z$), then $\psi(-1)=0$ as $k$ is odd, so $M^k-1$ and $M^k+1$ (both $\equiv\pm1\pmod M$)
are zeros of $\chi$, and $\chi(M^k)=k\,\chi(M)=0$; thus $(M^k-1,M^k,M^k+1)$ is a zero triple.

%======================================================================
\section{Data and reproduction}
All files are in the directory of this preprint.
\begin{itemize}
\item Lean files \path{lean/Main.lean} and \path{lean/Main7.lean}, with outputs
\path{lean/lean_output.txt} and \path{lean/lean_output_k7.txt}; control files
\path{lean/NegOpt.lean} and \path{lean/PosOpt.lean}. Check with
\texttt{lean Main.lean} run inside \path{lean/} (Lean 4.34.1, core library only; no \texttt{lakefile}
is needed).
\item \path{data/assignment_k5.txt}, \path{data/assignment_k7.txt}: the exceptional primes and their
values (all other primes have value $1$).
\item \path{scripts/}: C sources of the local search (\path{ws2.c}), the exact search (\path{lam.c}),
the two independent checkers (\path{vchk.c}, \path{segchk.c}), the $\Omega$-baseline (\path{om.c}),
and \path{reproduce.sh}.
\end{itemize}

{\small
\bibliographystyle{plain}
\bibliography{references}
}
\end{document}
